What are solutions to the bankrun problem and explain the simpler ones

- Suspension of convertibility
- Narrow Banking
- Deposit insurance
- Lender of Last Resort (LLR)

Suspension of convertibility
No more withdrawals allowed after some threshold fraction \( \hat{f} \) of
withdrawals
- Suspension need not be applied if announced reliably!
- Type 2 depositors won't run because they then "know" that bank has sufficient funds to repay them.

Solution not practical if fraction of early consumers \( \alpha \) is not
known upfront (e.g. 𝛼 is stochastic), ironically because if there are then more type 2
investors than expected there is not enough money allocated for them (as this smaller amount was supposed to go to type 1).

Narrow Banking
Restrict (risk and) maturity transformation by banks
Banks only allowed to invest into riskless and highly liquid money market instruments
=> Decreases GDP growth


Explain the pro's and cons of deposit insurance

Fraction of deposits is insured by some deposit insurance fund, which is typically Government financed,
since the Government as taxation authority, and is therefore less likely to default.
Alternatively a private deposit insurance fund, which is financed with annual fees by banks is possible.

This introduces Moral hazard:
Depositors have no incentive to monitor riskiness of bank activities and banks could increase risk without 
paying depositors higher interest rates.

In Germany Deposit insurance is required by law since 1998 with min. of 100k.

Regarding the private fund solution, should Banks pay a flat fee? 
No, they it should be based on risk.
Should depositors be provided with full insurance or some
deductible? Deductible, since full insurance is unreachable.


The incentives are in regard to moral hazards are modelled as follows:

Recall that we can simplistically model the bank (or a company generally) via its balance sheet 
at \( t = 0 \) as 
Aktiva 
\( S_0 \)
Passiva
Own capital
\( C(K, T) \)
Debt
\( \frac{K}{(1 + r)^T} - P(K, T) \)

and at \( t = T \) (maturity of debt) as
Aktiva 
\( S_T \)
Passiva
Own capital 
\( \text{max}\{S_T - K, 0\} \)
Debt
\( K - \text{max}\{K - S_T, 0\} \)

Therefore the equity holders of a bank hold a call option
if we consider the the debt \( K \) as the promised return to depositors,
and a Deposit insurance fund behaves like a put option, both with strike \( K \).


What is a Lender of Last Resort and how is it involved in bank runs


Typically, Central Banks act as lender of last resort:
They provide liquidity to illiquid, but otherwise solvent banks
(at least in theory)
Problems
distinguish between insolvent banks and illiquid banks
Moral hazard issue
price stability
exchange rate stability
(run on bank and currency Æ often ineffective)
asset price bubbles

The classical LOLR theory is Bagehot (1873):
lend early and freely (i.e., without limit) to solvent firms
against good collateral (collateral which in normal times is
considered good collateral)
at high rates

anks borrow from LOLR to substitute funding losses from
bank runs (or, more generally, when liquidity problems arise)
Changes only on Equity and Liability side of the balance sheet
Total assets stay the same!
Aimed to help illiquid, but solvent banks.
Ideally: Just announcing a LOLR solution should suffice to avoid bank runs.

In the ECB and often interbank market case, 
such business is done via Sale-and-repurchase agreement (repo)
E.g.
ECB buys asset (the collateral) with market value \( $100 \) for \( $80 \)
Bank agrees to repurchase next day for \( $88 \).
The repo rate is then \( \frac{88 - 80}{80} = 10\% \)
The repo rate is then \( \frac{100 - 80}{100} = 20\% \)
Haircut does not depend on borrower, but on collateral.


What is the Risk taking incentive theory / Risk shifting for Lenders of Last Resort?

Banks close to default want to take higher risk because Equity holders have no down side risk, 
but upside potential increases with risk taking
therefore they buy additional risky securities, i.e. extend total assets of the balance sheet.
Typically those assets reap high interest rates/returns on the risk without equity holder bearing the down side risk
and default also when the bank defaults (which creates even more risk due to correlation).
